%!TEX root = ../report.tex
\documentclass[report.tex]{subfiles}
\begin{document}

\chapter{Introduction}
\label{ch:intro}

\section{Motivation}
\label{sec:motivation}

A drone that has to act on what it sees needs a semantic label for every frame
it captures. Obstacle avoidance, landing-site selection and terrain following
all read from a dense map of the scene, and they read from it while the
aircraft is still moving.

Producing that map one frame at a time is expensive. The strongest
segmentation model trained for this project takes about \LatTypical~ms on a
native \SI{4}{K} frame, so labelling a \SI{45}{s} clip costs
\ClipMinutes~minutes of GPU time, roughly \RealtimeFactor$\times$ longer than
the flight it describes, and on an NVIDIA A100 rather than on anything a
drone could carry. Chapter~\ref{ch:results} reports the model and the
measurement in full.

Much of that work is also repeated. At \SI{20}{FPS} consecutive frames differ
by a few tens of pixels of apparent motion, so the network recomputes almost
the same mask twenty times a second. If the mask from one frame could be
carried forward to the next, the model would need to run only when carrying it
forward stops working.

That is the system evaluated here: segment one frame, move its mask onto the
frames that follow along optical flow, and re-run the model only when the moved
mask can no longer be trusted. Optical flow supplies the correspondence, since
a flow field states where each pixel went and a label attached to that pixel
can travel with it, and learned estimators such as RAFT~\cite{raft2020} and
SEA-RAFT~\cite{searaft2024} are accurate enough for this to be done directly.
The experiments use UAVid~\cite{uavid2020}, an urban-drone benchmark of
oblique, low-altitude video.

Carrying labels along flow is not itself new; what differs here is where the
first mask comes from. SegProp~\cite{segprop2020} and the video propagation of
Zhu et~al.~\cite{zhu2019} propagate \emph{human-drawn} annotations offline, to
grow a training set from a few labelled frames. In flight there is no
annotation to propagate, so the mask must be a model prediction produced on
the fly. That makes producing it the cost worth avoiding, and it forces the
system to judge for itself when the propagated mask has decayed enough to be
worth replacing.

\section{Challenges}
\label{sec:challenges}

\paragraph{Flow is the fixed cost.} RAFT-family
models~\cite{raft2020,searaft2024} build an all-pairs correlation volume whose
size grows \emph{quadratically} with resolution, so at \SI{4}{K} their memory
is dominated by that activation rather than by weights. More importantly, flow
must be computed for every consecutive pair before the staleness test can be
evaluated, so it is paid whether the model runs or not. This bounds the
achievable speedup.

\paragraph{Detecting bad flow is harder than computing it.} Flow fails at
occlusion boundaries, on fast-moving objects and under parallax, and the
propagated label is then wrong. The standard detector is the forward-backward
consistency check~\cite{sundaram2010}, which rejects a pixel whose round trip
through both flow directions does not return to its origin. It tests whether
the two flow fields \emph{agree}, not whether the label is \emph{correct}.

\paragraph{Ground truth is sparse.} UAVid labels one frame in a
hundred~\cite{uavid2020}, so every quality figure here rests on \NumGT{}
evaluation points across \NumSeq{} sequences: enough for a paired comparison of
two methods on identical frames, not enough to characterise the distribution of
propagation quality.

\section{Problem Statement}
\label{sec:problem}

Given a UAV video $\{I_0, \dots, I_{T}\}$, a segmentation model $f$ costing
$t_{\text{model}}$ per frame and a flow model costing $t_{\text{flow}}$ per
frame pair, produce a dense mask for \emph{every} frame in less than
$T \cdot t_{\text{model}}$ wall-clock time, with mean quality as close as
possible to running $f$ on every frame.

The policy evaluated is: run $f$ on frame $0$; propagate its output along flow;
maintain a cumulative estimate $\gamma$ of the fraction of the frame still
trustworthy; and re-invoke $f$ when $\gamma$ falls to a threshold $\tau$. Three
questions are asked of it: how much time it saves, how much quality that costs,
and what $\tau$ controls.

\section{Contributions and Scope}
\label{sec:contributions}

\begin{enumerate}
    \item \textbf{An adaptive hybrid inference pipeline} that propagates model
          predictions along SEA-RAFT flow and re-invokes the model on a
          forward-backward coverage trigger, running \Speedup$\times$ faster
          than model-every-frame inference at \Retention\% of its mIoU on the
          UAVid validation split (Section~\ref{sec:valresults}).
    \item \textbf{A component study} selecting the segmentation backbone
          (Section~\ref{sec:backbones}) and the flow backend
          (Section~\ref{sec:backends}) by measurement.
    \item \textbf{A characterisation of the coverage threshold} as a quality
          control rather than a speed control, which follows from flow being
          the fixed cost (Section~\ref{sec:sweep}).
    \item \textbf{A reproducible implementation} with interchangeable backbones
          and flow backends, chunked streaming flow that keeps peak memory flat
          regardless of video length, and a resolution setting that makes
          \SI{4}{K} flow tractable (Chapter~\ref{ch:solution}).
\end{enumerate}

The scope is narrow. Evaluation uses the UAVid validation split only, since the
test labels are withheld~\cite{uavid2020}; the flow-backend comparison uses
Ruralscapes for the reason given in Section~\ref{sec:backends}. Each backbone
was trained once with one seed, so backbone differences support a ranking and
not a significance claim. Propagated masks are used unrefined, and there is no
deployment engineering: no quantisation, distillation or TensorRT export, and
all timings come from one A100 under bf16 PyTorch.

\end{document}
